%chktex-file 8
\section{Task 2}
\makeatletter\def\@currentlabel{Task 2}\makeatother\label{sec:task2}

\subsection{Overview}
\begin{frame}{Overview}
  \begin{block}{Prompt}
    Classify the degraded test set using two non-exclusive approaches:
    unsupervised image enhancement at inference time, and
    augmentation-aware retraining of the classifier.
  \end{block}
  \vspace{\baselineskip}
  The degraded test set contains 500 images sharing the same ground truth as the clean test set,
  but acquired under strong degradation.
  Our workflow was to first measure the performance drop, then characterize the degradation
\end{frame}

\subsection{Baseline on the Degraded Set}
\begin{frame}{Baseline on the Degraded Set}
  We began by evaluating the Task 1 best model (\texttt{v1.12\_80\_epochs}, 83.6\% clean accuracy)
  directly on the degraded test set without any modification.
  \vspace{\baselineskip}
  \begin{columns}[onlytextwidth]
    \begin{column}{0.35\linewidth}
      \begin{itemize}
        \item [clean] \textit{Acc}: 83.6\%,\\ \textit{F1}: 0.827
        \item [degraded] \textit{Acc}: 56.2\%,\\ \textit{F1}: 0.554
      \end{itemize}
      \vspace{0.5\baselineskip}
      A \textbf{27-point drop}: the model had never seen degraded images during training
      and could not generalise to them.
    \end{column}
    \begin{column}{0.65\linewidth}
      \begin{center}
        \includegraphics[width=\textwidth]{../models/v1.12_80_epochs.png}
      \end{center}
    \end{column}
  \end{columns}
\end{frame}

\subsection{Approach A \textendash\ Unsupervised Enhancement}
\begin{frame}{Unsupervised Enhancement \textendash\ Techniques}
  Before modifying training, we attempted to correct the degradation \textit{at inference time}
  using classic unsupervised techniques:
  \vspace{0.5\baselineskip}
  \begin{itemize}
    \item \textbf{Auto Gamma Correction}: estimates a per-image gamma exponent from the mean
          brightness to pull it towards midgray (0.5).
          Dark images are brightened ($\gamma < 1$), bright images darkened ($\gamma > 1$).
    \item \textbf{White Balance --- Gray World}: assumes the average scene color is neutral grey;
          scales each RGB channel by $0.5\,/\,\mu_c$.
    \item \textbf{Contrast Stretch}: expands the dynamic range using the 8\textsuperscript{th}
          percentile as the black point, clipping and rescaling to $[0, 1]$.
  \end{itemize}
\end{frame}

\begin{frame}{Unsupervised Enhancement \textendash\ Results}
  \begin{center}
    \begin{tabular}{lcc}
      \hline
      Technique & Accuracy & F1 \\
      \hline
      No enhancement (baseline) & 56.2\% & 0.554 \\
      Auto Gamma Correction     & 47.6\% & 0.462 \\
      White Balance (Gray World) & 53.2\% & 0.530 \\
      WB + Gamma                & 45.0\% & 0.446 \\
      Contrast Stretch          & $\approx$52\% & --- \\
      \hline
    \end{tabular}
  \end{center}
  \vspace{0.5\baselineskip}
  Every technique worsened or negligibly changed the results.
  These methods all act on \textit{brightness}, \textit{contrast} and \textit{white balance}
  --- yet none of them improved things. \\
  \vspace{0.5\baselineskip}
  $\Rightarrow$ The degradation must be different in nature.
  We decided to understand what it actually is before going further.
\end{frame}

